\chapter{ВАЛИДАЦИЯ, ОГРАНИЧЕНИЯ И ОБСУЖДЕНИЕ РЕЗУЛЬТАТОВ}

% -------------------------------------------------------
\section{Статистическая валидация исследовательских гипотез}
% -------------------------------------------------------

\subsection{Гипотеза H1: универсальность предубеждений}

\textbf{Формулировка:} все тестируемые языковые модели демонстрируют статистически значимое предубеждение подтверждения и подверженность эвристике репрезентативности во всех трёх парадигмах.

\textbf{Вердикт: частично подтверждена.} Результаты по трём экспериментам согласованно указывают на широкую распространённость когнитивных предубеждений, однако одна модель (\texttt{qwen3.5-397b-a17b}) демонстрирует устойчивость к конъюнктивной ошибке.

В эксперименте 1 токен-bias в форме ошибки конъюнкции выявлен у восьми из девяти моделей. McNemar $p < 10^{-5}$ для семи моделей; delta\_bias варьируется от $+0{,}028$ (\texttt{glm-4.7-flash}) до $+0{,}565$ (\texttt{qwen3-next-80b}). \texttt{qwen3.5-397b-a17b} является единственной моделью без значимого bias: McNemar $p = 0{,}0625$, delta\_bias $= -0{,}005$, Cohen's $d = -0{,}112$ (пренебрежимый эффект).

В эксперименте 2 предубеждение подтверждения (совокупный индекс MBI + CBR) воспроизведено у всех девяти моделей в нейтральном промпт-условии Formal Logic Neutral: от 2,8\% (\texttt{gpt-5.2}) до 53,6\% (\texttt{gemma-3-12b-it}).

В эксперименте 3 CTR $\geq 0{,}629$ у всех non-reasoning моделей в baseline-условии. Даже наименее «подтверждающая» модель (\texttt{grok-4.1-fast}) в 63\% случаев тестирует тройки, которые её же текущая гипотеза предсказывает как True. Нормативный байесовский агент предполагает CTR, близкий к нулю.

\textbf{Оговорка:} хотя \texttt{qwen3.5-397b-a17b} устойчива к конъюнктивной ошибке, она демонстрирует выраженное предубеждение подтверждения в задачах Уэйсона (EMA $= 0{,}764$) и 2-4-6 (CTR $= 0{,}850$). Это означает, что устойчивость к одному типу предубеждения не гарантирует устойчивости к другому.

\subsection{Гипотеза H2: масштабирование и внутрисемейный парадокс}

\textbf{Формулировка:} в среднем более крупные модели демонстрируют более низкий уровень когнитивных предубеждений, однако данная зависимость не является строго монотонной внутри модельных семейств.

\textbf{Вердикт: подтверждена.} На уровне общей тенденции зависимость наблюдается: топовые модели (\texttt{gpt-5.2}, \texttt{claude-sonnet-4.6}, \texttt{qwen3.5-397b-a17b}) систематически демонстрируют более низкий уровень предубеждений. Однако внутрисемейный анализ выявляет выраженные инверсии.

В семействе Gemma направление эффекта зависит от типа задачи. В эксперименте 1 (вероятностное суждение) \texttt{gemma-3-27b-it} демонстрирует несравнимо более высокий токен-bias: delta\_bias $= 0{,}498$, fallacy rate $= 58{,}7\%$ против delta\_bias $= 0{,}039$, fallacy rate $= 2{,}9\%$ у \texttt{gemma-3-12b-it}. В эксперименте 2 (дедукция) картина инвертируется: EMA $= 0{,}578$ у 27b против $0{,}106$ у 12b. В эксперименте 3 (индукция) обе модели слабы (SR $= 0{,}04$ и $0{,}06$).

В семействе Qwen паттерн аналогичен: \texttt{qwen3-next-80b} при меньшем числе активных параметров демонстрирует существенно более высокий bias (delta\_bias $= 0{,}565$), чем \texttt{qwen3.5-397b-a17b} (delta\_bias $= -0{,}005$).

Наиболее правдоподобное объяснение: более крупная модель обладает более богатым пространством эмбеддингов, что позволяет ей лучше улавливать нарративную согласованность. В задаче Линды это парадоксально усиливает эвристику репрезентативности~--- модель увереннее «попадает в ловушку». В логической задаче Уэйсона улучшённая способность к распознаванию паттернов, напротив, помогает модели применять modus tollens. Таким образом, масштаб является неоднозначным фактором, эффект которого зависит от типа требуемого рассуждения.

\subsection{Гипотеза H3: спектральная природа предубеждений}

\textbf{Формулировка:} интенсивность когнитивных предубеждений систематически возрастает от вероятностной парадигмы к логической и достигает максимума в индуктивной.

\textbf{Вердикт: подтверждена.} Уровень предубеждений возрастает по спектру сложности:

\textit{Вероятностное суждение (эксперимент 1):} в uncorrelated-условии ошибок нет ни у одной модели (все $n_{12} = 0$). Предубеждение возникает исключительно под воздействием нарративной согласованности и полностью устраняется при семантически нейтральном хобби. Модели-лидеры (\texttt{qwen3.5-397b-a17b}, \texttt{glm-4.7-flash}, \texttt{gemma-3-12b-it}, \texttt{claude-sonnet-4.6}) демонстрируют fallacy rate $\leq 2{,}9\%$.

\textit{Дедуктивное рассуждение (эксперимент 2):} даже лучшая модель (\texttt{gpt-5.2}) не достигает EMA $= 1{,}0$ в нейтральном условии при нейтральном промпте ($0{,}884$). Слабые модели (\texttt{gemma-3-12b-it}: EMA $= 0{,}028$) практически не демонстрируют логической компетентности. Совокупный индекс bias (MBI $+$ CBR) достигает 53,6\% (от 0,4\% у \texttt{gpt-5.2} до 53,6\% у \texttt{gemma-3-12b-it}).

\textit{Индуктивное открытие правил (эксперимент 3):} CTR $> 0{,}60$ у всех non-reasoning моделей. Лучший Success Rate среди non-reasoning моделей~--- 28\% (\texttt{claude-sonnet-4.6}), что означает, что даже наиболее способная модель открывает правило менее чем в трети случаев. Переход к reasoning-режиму поднимает SR до 48\% (\texttt{gpt-5.2} reasoning + adaptive), однако это всё ещё ниже 50\%.

Таким образом, индуктивное открытие правил подтверждается как наиболее уязвимая область для современных LLM, что согласуется с теоретическим предсказанием о том, что многоходовое управление пространством гипотез требует вычислительных возможностей, отсутствующих в стандартной авторегрессионной генерации.

\subsection{Гипотеза H4: демографический токен-bias}

\textbf{Формулировка:} демографические переменные, кодируемые именем персонажа, оказывают статистически значимое влияние на интенсивность ошибки конъюнкции, при этом различия предположительно более выражены по оси расы, чем по осям возраста и пола.

\textbf{Вердикт: частично подтверждена.} Расовые эффекты присутствуют, но не являются универсальными. Статистически значимые расовые различия после BH-FDR коррекции обнаружены только у \texttt{glm-4.7-flash}: Asian ($p_{\text{BH-FDR}} = 0{,}0068$), Black ($p = 0{,}0137$), Hispanic ($p = 0{,}0182$), тогда как для White тест незначим. У остальных моделей McNemar $p < 10^{-15}$ для \textit{всех} расовых групп, что указывает на мощный и универсальный эффект конъюнктивного предубеждения, не зависящий от демографической подстановки.

Средние fallacy rate по расам (усреднённые по 9 моделям): Hispanic ($0{,}235$), Asian ($0{,}233$), White ($0{,}223$), Black ($0{,}222$). Разброс составляет менее 1{,}5 процентных пункта, что указывает на отсутствие \textit{универсального} расового смещения в конъюнктивной ошибке. Две модели (\texttt{gemma-3-12b-it}, \texttt{grok-4.1-fast}) демонстрируют реверсивный эффект: White-имена ассоциированы с несколько более высокой fallacy rate, чем minority-группы.

Возраст не дал значимого эффекта ни у одной модели: разница между возрастными группами с максимальной и минимальной fallacy rate составляет менее 0{,}3~п.п.

% -------------------------------------------------------
\section{Обсуждение результатов в контексте смежных исследований}
% -------------------------------------------------------

\subsection{Ошибка конъюнкции: сравнение с Suri et al.\ и Macmillan-Scott \& Musolesi}

Suri и соавторы~\cite{suri2024large} показали, что GPT-3.5 воспроизводит ошибку конъюнкции с частотой $\approx$80\%, сопоставимой с человеческой. Результаты настоящего исследования согласуются с этим выводом для части моделей: \texttt{qwen3-next-80b} демонстрирует fallacy rate $= 60{,}8\%$, а \texttt{gemma-3-27b-it}~--- $58{,}7\%$. Однако разброс между моделями значительно шире: от $0{,}2\%$ (\texttt{claude-sonnet-4.6}) до $60{,}8\%$, что не наблюдалось в работах, тестирующих одну-две модели.

Macmillan-Scott и Musolesi~\cite{macmillan2024irrationality} подтвердили конъюнктивную ошибку на нескольких моделях. Настоящее исследование дополняет их выводы двумя аспектами: (1) метрика Severity of Bias обнаруживает латентное предубеждение у \texttt{claude-sonnet-4.6} при единственной бинарной ошибке, что было бы невидимо в их методологии; (2) строго однонаправленная асимметрия ошибки ($n_{22} = 0$ и $n_{12} = 0$ для всех моделей) отличает паттерн LLM от человеческих данных, где часть испытуемых ошибается в обоих условиях~\cite{tversky1983extensional}.

Jiang и соавторы~\cite{jiang2024peek} доказали, что успешность LLM на логических задачах часто является артефактом токен-bias. Результаты эксперимента 1 согласуются с этим выводом: замена коррелированного хобби на некоррелированное полностью устраняет ошибку у всех моделей, что подтверждает определяющую роль нарративного контента, а не абстрактного понимания вероятности.

\subsection{Задача Уэйсона: контент-эффект и сравнение с Dasgupta et al.}

Dasgupta и соавторы~\cite{dasgupta2024language} обнаружили контент-эффекты на деконтаминированных стимулах: социальные контракты решаются точнее абстрактных правил. Результаты эксперимента 2 полностью подтверждают этот вывод (медианный прирост EMA $+0{,}084$) и углубляют его: условие «незнакомых фантастических социальных контрактов» даёт EMA, сопоставимое со знакомыми контрактами (медианный прирост от FL Neutral $+0{,}072$). Это поддерживает гипотезу прагматических схем Ченга и Холиоака~\cite{cheng1985pragmatic}: деонтическая структура обязательства фасилитирует логический вывод независимо от знакомости содержания.

Macmillan-Scott и Musolesi~\cite{macmillan2024irrationality} выявили высокую нестабильность ответов LLM при повторных запросах. Введение системы из пяти перестановок в настоящем исследовании количественно подтверждает это: Consistency Rate варьируется от 0{,}110 (\texttt{gemma-3-12b-it}) до 0{,}910 (\texttt{qwen3-next-80b}), что ставит под сомнение надёжность точечных single-shot оценок.

\subsection{Задача 2-4-6: сравнение с WILT и CogBench}

Banatt и соавторы~\cite{banatt2024wilt} в бенчмарке WILT показали, что все модели проявляют признаки предубеждения подтверждения. Результаты эксперимента 3 согласуются с этим выводом и конкретизируют его: CTR варьируется от 0{,}629 до 0{,}976, причём даже наилучший результат существенно выше нормативного уровня. Вместе с тем введение строгой метрики экстенсиональной эквивалентности (IOU на 200\,000 троек) обеспечивает более точную оценку, чем текстовое сравнение гипотез, использованное в WILT.

Batista и Griffiths~\cite{batista2026rational} связали предубеждение подтверждения в LLM с RLHF-выравниванием. Результат эксперимента 3 частично согласуется: adaptive-промпт, содержащий явные инструкции по фальсификации, снижает CTR у большинства моделей, однако это снижение не конвертируется в рост SR. Это указывает на то, что предубеждение подтверждения имеет как минимум два независимых компонента: выбор тестируемых троек (CTR) и качество ревизии гипотез (HCR),~--- и RLHF может влиять на них по-разному.

Бенчмарк CogBench~\cite{coda2024cogbench} зафиксировал существенный разрыв между LLM и нормативными агентами в задачах, требующих активного поиска фальсифицирующих свидетельств. Результаты эксперимента 3 подтверждают и углубляют этот вывод: даже reasoning-модели не достигают SR $= 0{,}50$, что означает, что ни одна из протестированных архитектур не способна к нормативному индуктивному рассуждению на данном классе задач.

% -------------------------------------------------------
\section{Практические рекомендации по митигации предубеждений}
% -------------------------------------------------------

Результаты трёх экспериментов позволяют сформулировать конкретные рекомендации для трёх групп потребителей.

\subsection{Для инженеров: выбор модели и стратегии промптинга}

\textbf{1. Тестировать CoT per-model.} Chain-of-Thought не является универсальной коррекцией предубеждений. В эксперименте 2 CoT улучшает EMA у всех трёх проприетарных моделей (медианный прирост $+0{,}180$), однако ухудшает результаты у двух из шести открытых моделей. Наиболее показателен случай \texttt{glm-4.7-flash}: CoT увеличивает Confirmation Bias Rate с $0{,}248$ до $0{,}548$. Рекомендация: перед развёртыванием CoT в продакшене протестировать его эффект на целевой модели с помощью предложенного фреймворка.

\textbf{2. Reasoning > prompt engineering для индуктивных задач.} В эксперименте 3 reasoning-режим обеспечивает качественный скачок (SR \texttt{gpt-5.2} reasoning + adaptive $= 0{,}48$), недостижимый через любые манипуляции промпта у non-reasoning моделей (максимум $0{,}28$). Для задач, требующих многоходового логического вывода, выбор модели с поддержкой reasoning более эффективен, чем оптимизация промптов.

\textbf{3. Контент задачи важнее демографии.} Демографические подстановки оказывают минимальное влияние на интенсивность конъюнктивной ошибки (разброс менее 1{,}5~п.п. между расовыми группами). В то же время профессиональная категория виньетки модулирует fallacy rate от 15\% (Finance) до 40\% (Social Services). При конструировании промптов для систем поддержки принятия решений следует фокусироваться на минимизации нарративной «гравитации» содержания, а не на демографической нейтрализации.

\subsection{Для исследователей: метрики и методология}

\textbf{1. Использовать непрерывные метрики вместо бинарных.} Метрика Severity of Bias обнаруживает предубеждение у \texttt{claude-sonnet-4.6} при единственной бинарной ошибке. Бинарная точность (Accuracy/Exact Match) способна пропустить латентное предубеждение в «потолочном» режиме.

\textbf{2. Оценивать стабильность через перестановки.} Consistency Rate выявляет модели с нестабильным рассуждением (\texttt{gemma-3-27b-it}: EMA $= 0{,}578$ при Consistency $= 0{,}280$), которые выглядят приемлемо при single-shot оценке.

\textbf{3. Использовать экстенсиональную эквивалентность для индуктивных задач.} Текстовое сравнение (exact match) неприменимо из-за синтаксической вариативности формулировок LLM. IOU на Монте-Карло-выборке обеспечивает математически строгую оценку.

% -------------------------------------------------------
\section{Ограничения исследования}
% -------------------------------------------------------

\subsection{Методологические ограничения}

\textbf{Контаминация обучающих данных.} Несмотря на принятые меры по генерации оригинальных датасетов (факториальный дизайн виньеток, нейтральные абстрактные правила, фантастические социальные контракты), полная гарантия деконтаминации недостижима: нельзя исключить присутствие структурно аналогичных стимулов в обучающих корпусах тестируемых моделей. Настоящее исследование измеряет итоговое поведение модели вне зависимости от механизма его порождения (memorization vs.\ emergent bias), однако разграничение этих механизмов остаётся открытым вопросом.

\textbf{Недетерминированность эксперимента 3.} Монте-Карло-оценка IOU не является полностью детерминированной при \texttt{temperature=0} на стороне модели из-за отсутствия фиксации numpy seed. Это означает, что точные численные значения SR и AUC-IOU могут незначительно варьироваться при повторных запусках.

\textbf{Ограниченность диапазона моделей.} Девять тестируемых моделей представляют актуальный срез на период проведения экспериментов (начало 2026~г.), однако не охватывают весь спектр существующих архитектур. За рамками исследования остались: специализированные модели для математического рассуждения, мультимодальные системы, модели с размером менее 7B параметров (за исключением \texttt{glm-4.7-flash}).

\textbf{Однократность прогонов.} Большинство экспериментальных условий тестировались в одном прогоне на модель. Для эксперимента 2 вариабельность частично контролируется через систему из пяти перестановок, однако для экспериментов 1 и 3 оценка внутрипрогонной дисперсии не проводилась.

\textbf{Самооценка уверенности как прокси.} Метрика Severity of Bias основана на запрошенной у модели численной оценке уверенности. Калибровка этих оценок варьируется между моделями и не верифицировалась независимым способом. Значения уверенности следует интерпретировать как относительный показатель в рамках одной модели, а не как абсолютно сопоставимую межмодельную величину.

\subsection{Ограничения отдельных экспериментов}

В эксперименте 1 демографический анализ был проведён только для моделей \texttt{glm-4.7-flash} и \texttt{qwen3.5-397b-a17b} в полном факториальном дизайне; для остальных моделей сетка была сокращена после обнаружения незначимости возрастного и гендерного эффектов. Это решение, хотя и обоснованное промежуточными результатами, ограничивает возможность сравнительного анализа расового эффекта по всем девяти моделям.

В эксперименте 2 все три промпт-условия (Original, Neutral, CoT) оценивались только на датасете Formal Logic Neutral, что не позволяет оценить взаимодействие промпт-условий с контентными условиями.

В эксперименте 3 reasoning-режим тестировался только для двух моделей (\texttt{gpt-5.2} и \texttt{claude-sonnet-4.6}), что ограничивает выводы о его универсальности.

% -------------------------------------------------------
\section{Апробация результатов}
% -------------------------------------------------------

Основные положения настоящего исследования были представлены на кандидатском экзамене в AI Talent Hub ИТМО в январе 2026~г. В ходе представления были получены замечания от научного руководителя и членов комиссии, часть которых учтена при доработке методологии эксперимента 3 (введение условия ранней остановки по порогу IOU~$>~0{,}95$) и формулировке ограничений исследования.

По результатам проведённых экспериментов подготовлена рукопись статьи, планируемая к подаче в рецензируемое издание в области искусственного интеллекта и когнитивных наук. Экспериментальный код и датасеты публикуются в открытом доступе для обеспечения воспроизводимости результатов.

% -------------------------------------------------------
\section{Направления будущих исследований}
% -------------------------------------------------------

Настоящее исследование обозначает ряд перспективных направлений для дальнейшей работы.

\textbf{Расширение набора парадигм.} Три использованные парадигмы охватывают вероятностный, дедуктивный и индуктивный типы рассуждения, однако спектр когнитивных предубеждений значительно шире. Перспективными для тестирования LLM являются: эффект якорения, пренебрежение базовыми вероятностями (base-rate neglect), эффект фрейминга в контексте потерь и приобретений, а также предубеждения, специфичные для причинно-следственного рассуждения.

\textbf{Факторный дизайн reasoning $\times$ adaptive $\times$ парадигма.} Reasoning-режим тестировался только в эксперименте 3 и только для двух моделей. Систематическое сравнение reasoning и non-reasoning режимов на всех трёх парадигмах позволит определить, является ли качественный скачок в индуктивном рассуждении специфическим свойством задач многоходового вывода или универсальным эффектом.

\textbf{Исследование внутрисемейного парадокса.} Инверсия направления эффекта масштабирования между вероятностной и логической задачами в семействе Gemma требует механистического объяснения. Анализ голов внимания на демографические и нарративные токены у 12b и 27b версий модели позволит определить, связан ли парадокс с различиями в структуре эмбеддингов или с post-training выравниванием.

\textbf{Механизм отрицательного CoT-эффекта.} Две модели (\texttt{glm-4.7-flash}, \texttt{gemma-3-27b-it}) демонстрируют ухудшение результатов при CoT-промпте. Исследование причин этого эффекта --- связано ли оно с длиной контекста, отвлечением внимания на рассуждение или с принципиальной неспособностью модели генерировать корректные промежуточные шаги --- имеет самостоятельную научную ценность.

\textbf{Динамика предубеждений при файнтюнинге.} Открытым остаётся вопрос о том, как целенаправленное обучение на нормативно корректных примерах рассуждения влияет на профиль когнитивных предубеждений. Исследование эффектов RLHF и supervised fine-tuning на метриках настоящего фреймворка представляет как теоретический, так и прикладной интерес.

\textbf{Детерминированность и воспроизводимость метрики IOU.} Добавление фиксации numpy seed в функцию \texttt{compare\_sets} и верификация стабильности результатов на нескольких независимых прогонах является первоочередной технической задачей для повышения воспроизводимости эксперимента 3.

\textbf{Калибровка метрики Severity of Bias.} Верификация калибровки запрошенных оценок уверенности через независимые методы (например, температурное масштабирование или сравнение с логарифмическими вероятностями) позволит повысить межмодельную сопоставимость метрики.

\textbf{Мультимодальные модели.} Все тестируемые модели являются текстовыми. Мультимодальные системы, обрабатывающие изображения совместно с текстом, могут демонстрировать качественно иные профили предубеждений при решении задач, допускающих визуальное представление.
